\section{模型的评价}

\subsection{模型的优点}

\begin{itemize}[itemsep=2pt, topsep=3pt]
	\item \textbf{物理内核统一、四问口径一致}。四个问题共用式~\eqref{eq:balance}$\sim$\eqref{eq:curtail} 一组约束，各问之间只改变信息集与费用项，从根本上避免了“各问口径漂移”这一常见失分点。$17$ 套年度方案的最大母线平衡残差低于 $10^{-12}$ kWh。
	\item \textbf{信息集边界严格且可自动验证}。所有预测器只使用决策时刻之前已发生的观测，模型不依赖任何未来真实数据；并进一步用四组 mutation 测试把“无未来信息泄漏”从人工审查变成了可重复执行的自动化验收项。
	\item \textbf{风险定价有理论依据}。本文没有凭经验设定安全边际，而是从 $5$ 倍紧急电价推导出 $4:1$ 的非对称损失，再由报童模型证明最优光伏点预测应取预测分布的 $0.2$ 分位数，并通过独立数值实验验证其最优性。该结论使“保守预测”从直觉上升为可证明的最优决策。
	\item \textbf{状态递推贯穿始终}。储电量 $S_{d,t}$ 作为唯一的跨时段状态变量，在每一天、每一个预报时刻都以上一时段\emph{实际测得}的储电量为初值重新优化，因此各期方案互不相同，符合状态递推类调度问题（如 2021 年 C 题周—周库存递推）的建模要求。
	\item \textbf{结论可归因}。所有对比实验都只改变一个因素：问题二的 A/B/C/D 分解了完美负载信息、光伏风险定价与实时储能反馈三个来源；问题三的 $S_0\sim S_3$ 消融把“是否需要引入新增预报”转化为可量化的增量价值；问题四的 perfect/causal 分离出电价信息的价值上限。
	\item \textbf{结算口径处理严谨}。两种结算解释分别建立目标函数并\textbf{独立重新优化}，而不是对同一调度事后换算式计费；分段费用函数通过引入 $u,v$ 并施加等式约束精确线性化，线性化误差小于 $10^{-12}$ 元。
	\item \textbf{工程可复现}。四问均有独立求解脚本，输出保留完整的环境信息、随机种子与结果哈希；legacy 参数组合下可无损复现修订前的全部数值。
\end{itemize}

\subsection{模型的不足}

\begin{itemize}[itemsep=2pt, topsep=3pt]
	\item 问题二与问题三采用“计划终端储电量等于当日实际初值”的日循环口径。该处理能抑制日末透支、保证储能长期可持续，但限制了跨日峰谷套利。灵敏度分析表明该约束对总费用的影响不足千分之二，但在更长的规划周期中其代价会上升。
	\item 光伏预测未引入数值天气预报、辐照度、温度与云量等外生变量，仅使用同刻历史序列与残差分位数。当气象形态发生突变（例如连续阴雨转晴）时，预测与安全边际的调整会滞后一天左右。
	\item 储能模型忽略自放电、容量衰减、循环寿命成本与效率随功率变化的特性，因此年度费用中没有反映电池的长期老化代价，得到的充放电深度可能比实际最优更大。
	\item 问题四的 perfect price 情形假设 $0{:}00$ 即可获知当天全部 $144$ 个电价。这在多数电力市场并不成立，本文已用 causal 情形给出可实现结果，但价格预测本身的精度仍有提升空间。
	\item 问题二采用逐日滚动的确定性等价形式求解两阶段模型，安全边际由历史残差分位数估计。当可用历史不足（如年初的前 $7$ 天）时，分位数估计的稳定性有限，本文以附件 1 的典型日曲线作为冷启动先验，但这仍是一个近似。
\end{itemize}

\section{模型的改进与推广}

\subsection{模型的改进}
\label{sec:improve}

\textbf{（1）把单日滚动扩展为多日滚动时域}。当前模型每天重新生成基准计划，跨日只传递储电量。若把优化窗口扩展到 $3\sim7$ 天并只执行首日决策，同时在窗口末端加入基于未来价格与负载的储能残值函数，就可以量化跨日套利收益并进一步降低边界效应。

\textbf{（1$'$）用混合整数规划消除同时充放的线性松弛退化}。如 \ref{sec:validation} 节所披露，线性规划无法表达储能的互补条件 $C_tH_t=0$，在 additive 口径下会出现“边充边放”的内部耗散伪解。引入 $0$-$1$ 变量 $z_t$ 并施加 $C_t\le \bar E z_t$、$H_t\le \bar E(1-z_t)$ 即可精确刻画，代价是滚动优化由线性规划升级为混合整数线性规划。以本文 $144$ 时段窗口、$144$ 个 $0$-$1$ 变量的规模估计，单次求解仍在百毫秒量级，属于低成本、高收益的改进。

\textbf{（2）把确定性等价升级为显式随机规划}。本文已证明 $5$ 倍电价对应 $4:1$ 非对称损失，并据此得到 $0.2$ 分位数的解析最优解。若把负载与光伏的联合分布显式建模，可进一步构造基于情景抽样（SAA）或条件风险价值（CVaR）的两阶段随机规划，在保留分位数结论的同时给出风险度量。

\textbf{（3）引入更完整的价格预测}。可在电价预测中加入日内周期项、工作日／节假日效应与温度特征，或采用分位数回归直接给出价格分布，从而把电价不确定性也纳入目标函数而非仅做信息集分层。

\textbf{（4）细化储能模型}。可把循环深度相关的寿命损耗、效率曲线、自放电与维护成本写入目标函数，并用历史运行数据重新标定参数；当模型规模增大时，可采用模型预测控制与 Benders 分解等算法保持滚动求解速度。

\subsection{模型的推广}

本文提出的“\emph{统一物理内核 $+$ 信息集分层 $+$ 文献可推断的风险定价}”框架并不依赖于特定题目，可直接迁移到下列场景：

\begin{itemize}[itemsep=2pt, topsep=3pt]
	\item \textbf{园区综合能源系统}：在电、热、气多能耦合场景下，把母线能量平衡替换为多能流平衡，把储能状态递推扩展到储热、储气装置，信息集分层与滚动结构保持不变。这类多能耦合微网的日前经济调度已有成熟建模范式可供对接 \cite{liu2018dayahead}。
	\item \textbf{电动汽车充换电站}：把负载替换为充电需求、把储能替换为动力电池，即可用同一框架优化峰谷套利与需求响应。
	\item \textbf{工业用户侧储能}：在需量电费与分时电价并存的场景下，只需在目标函数中增加最大需量项，其余约束与信息集处理方式完全复用。
	\item \textbf{含风光互补的微网}：把光伏预测替换为风光联合预测，按各电源的预测误差分别设定非对称损失权重与分位数水平即可。
	\item \textbf{更一般的状态递推型调度问题}：凡是“每期决策依赖上一期实际状态、且计划与执行存在偏差结算”的问题（如库存—订购—运输、水库调度、供应链补货），都可以套用本文的“状态递推 $+$ 单因素消融 $+$ 信息集自动验证”方法。
\end{itemize}
